\begin{table}[htbp]
\centering
\caption{Covariate Balance: 1940 State Characteristics on Mobilization}
\label{tab:balance}
\small
\begin{tabular}{lccc}
\hline\hline
1940 Covariate & Coefficient & Std. Error & $p$-value \\
\hline
Pct urban & 0.0600 & (0.0375) & 0.116 \\
Pct black & -0.0620*** & (0.0213) & 0.006 \\
Pct farm & -0.0650** & (0.0281) & 0.025 \\
Mean educ & 0.6047*** & (0.1553) & 0.000 \\
Mean age & 0.4226*** & (0.1495) & 0.007 \\
Pct married & -0.0081 & (0.0059) & 0.178 \\
Pct female & -0.0007 & (0.0020) & 0.724 \\
\hline
Joint $F$-test & \multicolumn{2}{c}{$F = 2.60$} & 0.031 \\
\hline\hline
\end{tabular}
\begin{tablenotes}
\small
\item \textit{Notes:} Each row reports the coefficient from a bivariate regression of the 1940 state characteristic on standardized mobilization rate, weighted by 1940 female population. The joint $F$-test regresses mobilization on all covariates simultaneously. Significant coefficients indicate mobilization is not conditionally random, motivating the inclusion of controls in main specifications.
\end{tablenotes}
\end{table}
